% Transcription — fonds Alexandre Grothendieck, shelfmark 52, batch 1 (pages 1-13).
% Established from the facsimile archives/batches/52/batch-01.pdf (a mirror of
% https://grothendieck.umontpellier.fr/52.pdf), per the skill
% transcribe-grothendieck. The folder is thirteen pages and this is its only
% batch. Page 1 is a tan cover sheet carrying two words of his in the top right
% corner and nothing else, "Th. de Torelli"; it therefore takes no \page marker
% and its words open the file as the section heading. The inventory's title
% "Théorème de Torelli" is that abbreviation expanded, so the title is his and
% not the archivists'. Pages 2 to 13 are the run itself and every one of them
% is mathematics; there is no unrelated verso anywhere in the folder, and no
% pagination of his own.
% The run is one argument in three movements. Pages 2 to 4 are the injectivity
% of Aut_S(C) -> Aut_{S-gr}(J^0_{C/S}) for a family of curves of genus >= 2,
% with the Lefschetz fixed-point count -1+2g-1 = 2(g-1) > 0 as its whole proof,
% and its three corollaries, the third being that Isom_S(C,C') -> Isom_S of the
% jacobians is a closed immersion. Pages 4 to 7 set up the moduli problem this
% is meant to serve: A a polarised abelian scheme over S, F(T) the isomorphism
% classes of curves over T with an isomorphism of their jacobian onto A_T, F'
% its rigidified variant, and the theorem that F' is representable by an M'
% quasi-projective over S on which A acts, M = M'/A. Pages 8 and 9 are the
% infinitesimal statement — a first-order deformation of C inside A is induced
% by an infinitesimal translation of A — read off the normal sequence
% 0 -> T -> V -> N -> 0 and reduced by duality to the surjectivity of
% H^0(C,Omega) (x) H^0(C,Omega) -> H^0(C,Omega^{(x)2}), which holds exactly when
% C is not hyperelliptic. Pages 10 to 13 bring in Matsusaka's criterion — an
% abelian variety carrying an irreducible positive divisor Theta with
% deg(Theta^n)/n! = 1 and Theta^{n-1} = (n-1)! Z is a jacobian — and close on
% three corollaries about descent of the situation to a base S.
% The hand is drafting throughout and this is the folder's governing fact.
% Following folder 29 batches 8 and 9 and folder 49, the statements, the
% displayed formulas and the underlined words are transcribed in full and the
% connective prose is left \ill{} wherever the leaf will not support a reading.
% Pages 3, 8 and 9 are written fair and read almost end to end; pages 5, 6, 7,
% 10 and 12 are a fast palimpsest, several of them carrying blocks cancelled by
% long diagonals over the top of an already crowded sheet, and much of their
% prose is not recoverable at the resolution the film carries. The apparatus is
% dense in consequence, and that density is the record of what these leaves
% actually support.
% Three things about the notation are fixed once here rather than page by page.
% He doubles the underline of a functor or a representing scheme — Pic, Aut,
% Isom, Hilb, and the sheaves T, V, N, Omega — and following folders 47 and 50
% the double rule is set as a single \underline. He writes the jacobian with an
% ornate capital J, J^0 for Pic^0 and J^1 for the torsor of degree 1 in which C
% sits, and the two exponents are distinguished on the page. And he writes
% Omega with a caron for the dual of Omega, i.e. the tangent sheaf; that is
% transcribed \check{\Omega} where the caron is on the leaf.
% Pass: Opus 5 (claude-opus-5), 2026-09-13 — first pass, unchecked against the pages by a human.
\documentclass[11pt,a4paper]{article}
% Le préambule commun à toutes les transcriptions du fonds Grothendieck.
%
% Il tient en une idée : **l'incertitude se compose, elle ne se résout pas.**
% Une transcription qui lisse un mot douteux a détruit la seule chose pour
% laquelle on la faisait — un lecteur doit pouvoir savoir, à chaque endroit,
% ce qui était sur le feuillet et ce qui a été ajouté. Les six macros qui
% suivent sont donc l'appareil critique, et non de la décoration.
%
% Les mêmes commandes sont reconnues par `scripts/render.mjs`, qui produit la
% vue de lecture du site. Une macro ajoutée ici doit l'être aussi là-bas,
% faute de quoi le rendu échouera bruyamment — ce qui est voulu : un rendu
% silencieusement faux est pire qu'un rendu absent.

\ProvidesPackage{grothendieck}[2026/08/08 Transcriptions du fonds Grothendieck]

\usepackage{amsmath,amssymb,amsthm}
\usepackage{mathrsfs}
\usepackage{stmaryrd}
% Les diagrammes commutatifs sont partout dans ce fonds : c'est la langue
% même de la géométrie algébrique de Grothendieck, et une transcription qui
% les rendrait en prose ne transcrirait rien.
\usepackage{tikz-cd}
% Français par défaut — c'est la langue des feuillets — mais l'anglais est
% chargé aussi : la traduction et le résumé sont des documents anglais, et la
% typographie française (espaces avant les deux-points) y serait une faute.
% Ces documents basculent par \selectlanguage{english} en tête de corps.
\usepackage[english,french]{babel}
\usepackage{fontspec}
\usepackage{geometry}
\geometry{a4paper,margin=25mm,marginparwidth=28mm}
\usepackage{marginnote}
\usepackage{xcolor}
% ulem, not soul. `soul`'s \st and \ul are built on a character-by-character
% scan that XeLaTeX's font machinery breaks: the document fails with an
% undefined control sequence rather than degrading. ulem does the same two jobs
% and is Unicode-engine safe. `normalem` stops it hijacking \emph, which the
% transcriptions use for Grothendieck's own emphasis.
\usepackage[normalem]{ulem}
\usepackage{hyperref}
\hypersetup{colorlinks=true,linkcolor=black,urlcolor=black,pdfborder={0 0 0}}

% --- Métadonnées du lot -----------------------------------------------------
% Reprises telles quelles par le script de rendu, qui les affiche en tête de
% la vue de lecture. Elles fixent ce que le fichier prétend couvrir : sans
% elles, une transcription flotte sans point d'attache dans un fonds de seize
% mille feuillets.

\newcommand{\folder}[1]{\gdef\grothFolder{#1}}
\newcommand{\batch}[1]{\gdef\grothBatch{#1}}
\newcommand{\leaves}[2]{\gdef\grothFirst{#1}\gdef\grothLast{#2}}
\newcommand{\foldertitle}[1]{\gdef\grothFoldertitle{#1}}
\gdef\grothFolder{?}\gdef\grothBatch{?}\gdef\grothFirst{?}\gdef\grothLast{?}\gdef\grothFoldertitle{}

% --- Le résumé --------------------------------------------------------------
%
% Il ouvre la lecture modernisée, sous le titre « Résumé », et il est composé
% en petit. Ce n'est pas un ornement : c'est le seul passage du document qui
% ne soit pas des mathématiques, et un lecteur doit pouvoir le reconnaître
% avant de l'avoir lu — puis le sauter s'il connaît déjà le sujet, ou s'y
% arrêter s'il ne le connaît pas. Le filet qui le suit marque la reprise du
% corps.

\newenvironment{resume}
  {\small\setlength{\parskip}{0.45em}}
  {\par\medskip\hrule\medskip}

% --- Mots-clés ---------------------------------------------------------------
%
% En anglais, en clôture du résumé de la lecture modernisée : le vocabulaire
% moderne sous lequel chercher ce que les feuillets construisent. Le manifeste
% du site (`npm run manifest`) les extrait pour étiqueter la cote entière —
% cette ligne est la source unique des tags, il n'y a pas de fichier à côté.
\newcommand{\keywords}[1]{\par\smallskip\noindent
  {\footnotesize\textsc{Keywords} --- \textit{#1}}\par}

% --- L'appareil critique ----------------------------------------------------

% Le numéro de feuillet, en marge. C'est l'ancre qui permet de régler un
% désaccord sur une formule en regardant *un* feuillet plutôt qu'une chemise
% de six cents. Le site s'en sert aussi pour tourner les pages du fac-similé
% quand on fait défiler la transcription.
\newcommand{\leaf}[1]{%
  \marginnote{\footnotesize\color{gray!70}\textsf{#1}}%
  \label{leaf:#1}\ignorespaces}

% Illisible : on ne devine pas. Un mot inventé qui se lit comme les autres est
% le pire résultat possible.
\newcommand{\ill}{\textcolor{red!60!black}{[\,\ldots\,]}}

% Lecture douteuse : le mot est donné, et signalé comme incertain.
\newcommand{\uncertain}[1]{\uline{#1}}

% Ajout de l'éditeur — une lettre manquante, un mot sous-entendu.
\newcommand{\add}[1]{\textcolor{blue!50!black}{[#1]}}

% Biffé par Grothendieck lui-même. Ce qu'il a rayé fait partie du document :
% une rature dit souvent où la pensée a changé de direction.
\newcommand{\struck}[1]{\sout{#1}}

% Note du transcripteur, sur l'état matériel du feuillet.
\newcommand{\note}[1]{\par\smallskip{\small\color{gray!60!black}\textit{[#1]}}\par\smallskip}

% Note marginale de Grothendieck, distincte du corps du texte.
\newcommand{\marginal}[1]{\par\smallskip\noindent\hspace{1em}%
  {\small\color{gray!60!black}#1}\par\smallskip}

% --- Titre ------------------------------------------------------------------

\AtBeginDocument{%
  \begin{center}
    {\large\bfseries Fonds Alexandre Grothendieck --- cote n\textsuperscript{o}~\grothFolder}\\[2pt]
    {\itshape\grothFoldertitle}\\[4pt]
    {\small feuillets \grothFirst--\grothLast\ (lot \grothBatch)}\\[2pt]
    {\footnotesize\color{gray!60!black}Transcription établie d'après le fac-similé numérisé
     par l'Université de Montpellier.\\
     Les lectures incertaines sont soulignées, les illisibles notées [\,\ldots\,],
     les ajouts entre crochets.}
  \end{center}
  \vspace{1em}}

\endinput


\folder{52}
\batch{1}
\pages{1}{13}
\dating{s.d.}
\watermark{Édition de démonstration}
\foldertitle{Théorème de Torelli : notes manuscrites (s.d.)}

\begin{document}

\section*{Th. de Torelli}

\note{titre de sa main, écrit au coin supérieur droit d'un feuillet de couverture
qui ne porte rien d'autre. Le feuillet ne reçoit donc pas de \texttt{page},
conformément à l'usage suivi depuis le dossier 29 : le saut dans la numérotation
est la trace qu'il a été passé. Le titre de l'inventaire, « Théorème de Torelli »,
est cette abréviation développée, et il est donc de lui}

\page{2}

\textbf{Th. 1}\quad Soit $C$ un schéma en courbes simple, propre, irréductible
sur $S$ \marginal{loc. noeth.}, de genre $g \geq 2$, $J^{0}_{C/S} =
\underline{\mathrm{Pic}}^{0}_{C/S}$ son schéma jacobien
\uncertain{homogène}.\note{le dernier mot de l'énoncé n'est pas assuré : les
traits se laissent lire « homogène » et rien d'autre ne s'y ajuste, mais
l'adjectif surprend appliqué à $\mathrm{Pic}^{0}$ lui-même} Alors
$\underline{\mathrm{Aut}}_{S}(C) \to
\underline{\mathrm{Aut}}_{S\text{-gr}}(J^{0}_{C/S})$ est \emph{injectif}.

Soit $\sigma$ un automorphisme de $C$ induisant l'identité sur $J^{0}_{C/S}$.
Soit $\bar{\sigma}$ l'automorphisme de $J^{1}_{C/S}$ induit par $\sigma$, il
sera, donc, de la forme $T_{s}$, \uncertain{translation}\note{le mot est abrégé
et la lecture repose sur ce que dit la suite de la phrase, non sur l'encre : ce
qui est associé à la section $s$ ne peut être que la translation par $s$}
associé à une section $s$ de $J^{0}_{C/S}$ (l'énoncé \ill{} l'énoncé
ensembliste). Il faut montrer que la section $s$ est l'identité. À vrai dire,
$C$ étant projectif \struck{sur} $S$, l'énoncé est équivalent à

\textbf{Corollaire 1}\quad $\underline{\mathrm{Aut}}_{S}(C) \to
\underline{\mathrm{Aut}}_{S\text{-gr}}(J^{0}_{C/S})$ est un
\emph{monomorphisme} (et $=$ une \emph{immersion fermée}, \ill{} \ill{}
rappelle que $\underline{\mathrm{Aut}}_{S}(C) \to$ \ill{} projectif sur $S$)

\struck{\ill{}}

\note{deux lignes insérées ici et biffées d'un trait continu ; elles ne se
laissent pas lire}

Comme les deux \ill{} sont \uncertain{non ramifiés} sur $S$ [pas
d'automorphismes infinitésimaux !] \ill{} \ill{} \ill{}

\page{3}

\textbf{Lemme}\quad Soit $C$ une courbe alg. non sing. complète
\uncertain{et} irréd. de genre $g \geq 2$ sur $k$ alg.\ clos, $\sigma$ un
automorphisme de $C$ induisant l'identité sur $J^{0}_{C/S}$. Alors $\sigma$
est l'identité.

En effet, l'automorphisme de $J^{1}_{C/S}$ induit par $s$ est de la forme
$T_{a}$, où $a$ est un pt de $J^{0}$ rat.\ $/k$, égal à $\sigma x - x$ pour tt
$x \in C(k)$. ($C$ étant considéré comme \uncertain{plongé} dans $J^{1}$.) Donc
$\sigma(x) = x + a$, donc \ill{} \ill{} \uncertain{en suivant} \ill{}
\uncertain{$\mathrm{Spec}(k)$}.

\textbf{Corollaire 2}\quad \emph{Sous les conditions du Th.}, \struck{Toute}
Une section $s$ de $J^{0}_{C/S}$ sur $A$\note{la lettre est un $A$ majuscule et
non un $S$ ; aucun $A$ n'a été introduit à ce point du dossier, il ne le sera
qu'à la page 4} telle que la translation $T_{s}$ de $J^{1}_{C/S}$ conserve $C$
(plongée dans $J^{1}_{C/S}$) est \struck{l'identité} section
nulle.\note{« section nulle » en interligne au-dessus de « l'identité » biffé}

En effet, il résulte de la formule des traces que le « nb.\ de pts fixes » de
$\sigma$ est égal : \struck{\ill{}} $-1 + 2g - 1 = 2(g-1) > 0$ (puisque
$\sigma$ opère trivialement dans la jacobienne $J^{0}$, donc sa trace vaut
$2g$] et si $x$ est un tel pt, \ill{} $a = \sigma x - x = 0$.

\note{la version biffée du calcul, écrite sous celle qui reste, n'aboutit pas
au même nombre ; elle est trop chargée pour être lue}

\page{4}

\note{le haut de la page est un bloc encadré, puis annulé par deux longues
diagonales. Il se lit en partie et est donné ici ; l'annulation n'est signalée
qu'une fois}

\struck{Utilisant le corollaire 2 lorsque $S = \mathrm{Spec}(k(\varepsilon))$
pour $S' = \mathrm{Spec}(\underline{D}(\underline{I}))$ : $\underline{I}$
\ill{} faisceau cohérent sur $S$, $\underline{D}(\underline{I}) =
\underline{O}_{S} + \underline{I}$ [$\underline{I}$ de carré nul] et \ill{}
\ill{} \ill{} $C' = C \times_{S} \ill{}$ qui induisent l'identité sur \ill{},
on trouve}

\struck{Prop.}

\textbf{Corollaire 3}\quad Soient $C$, $C'$ deux courbes sur $S$ satisfaisant
les hyp.\ du th. Alors
\[
  \underline{\mathrm{Isom}}_{S}(C, C') \longrightarrow
  \underline{\mathrm{Isom}}_{S}(J^{0}_{C/S}, J^{0}_{C'/S})
\]
est une \emph{immersion fermée}.

Soit maintenant $A$ un \emph{schéma abélien} \add{de dim.\ relative $g$}
\emph{polarisé} sur $S$ \ill{} au sens \emph{strict} : on se donne une section
du schéma de Nér.-Sév.). Pour tout $T$ sur $S$, on désigne par $F(T)$
l'ensemble des classes, à isom.\ près, de courbes $C$ sur $S_{T}$
\uncertain{$= T$}, satisfaisant aux conditions générales du \uncertain{Th.}, et
\emph{munies} d'un \emph{isom.} de $J^{0}_{C/T}$ sur $A$, \add{compatible avec}
\struck{\ill{}} les polarisations.

\page{5}

\note{la page est écrite très vite et à demi recouverte de biffures ; ce qui
suit est ce que le feuillet soutient}

$F'(T)$ l'ensemble analogue, mais où \ill{} \ill{} \ill{} \ill{} \ill{}
\ill{} \ill{} la donnée d'une rigidification par des éléments de
$(\underline{\Omega}^{1}_{C/T})^{\otimes h}$ ($h \geq 3$ \ill{}).

Soit $v = 2h(g-1) - (g-1) = (2h-1)(g-1)$.

\marginal{rédiger proprement}

\textbf{Théorème}\quad \struck{Le foncteur} \ill{} $F'(T)$ \struck{\ill{}}
est \emph{représentable} \ill{} par un $M'$ \ill{} quasi-projectif sur $S$,
\ill{} \ill{} \ill{}

\ill{} \ill{} $F'(T)$, \ill{} : \ill{} représentable \ill{} \ill{} \ill{}
application \ill{} \ill{} \ill{} \ill{} \struck{$F'(T)$} $\to F(T)$, le
passage \ill{}

\ill{} $F' \times_{F} F' \rightrightarrows F'$ et \ill{} d'ailleurs, \ill{}
$M'$ \ill{} \ill{} $M'$ \ill{} par des opérations de \ill{}

\note{le bord gauche de la page porte, en biais, un second bloc plus rapide
encore : on y lit « ex : \ill{} courbes \ill{} hyperelliptiques \ill{} », et
rien de plus ne s'en laisse tirer. Un petit cadre quadrillé, biffé, occupe le
bas de la marge}

\page{6}

\note{page de même farine que la précédente : les articulations portent, la
prose non}

\ill{} \ill{} \ill{} \ill{} par la théorie de \emph{Hilbert} [car le \ill{}
\ill{} \ill{}]. \ill{} \ill{} condition \ill{} \ill{} \ill{} plus petit
\ill{}

\ill{} \ill{} \ill{} \ill{} que $F'$, \struck{\ill{} $F$} \ill{} \ill{}
proportionnels, dans $F$ \ill{} \ill{} \ill{} \ill{} \ill{} l'\ill{} \ill{}
difficile être \ill{}. Pour le \ill{} \ill{} aux \struck{deux} cas
particuliers

(i)\quad $S = \mathrm{Spec}(k)$ \ill{} \ill{} \ill{}

\ill{} \add{deux} \ill{} \ill{} de \emph{Torelli} \ill{} sous sa forme
\emph{classique} \ill{}

\ill{} $M$ ne peut avoir \ill{} \ill{} \ill{} $Z/2 \in A_{-g}(A)$ \ill{} \ill{}
c'est \ill{} le \emph{théorème} \ill{}

\ill{} $M \to \mathrm{Spec}(k)$ \ill{} \ill{} \add{deux} \struck{le} cas de
$M$ \ill{} localisés \ill{} \ill{} identiques.

\page{7}

$F'$ \ill{} \ill{} le th.\ 1 \struck{\ill{}}, il opère sur \ill{} \ill{}
donc le \emph{quotient} \ill{} \ill{} \ill{}

\emph{si} $F'$ est représentable \ill{} par ($M'$, dans \ill{} $A$ qui
\ill{} $M'$, \ill{} \ill{} \ill{} le th.\ 1, il \ill{} \ill{} \ill{}. D'après
\ill{} \ill{} \ill{} $F$ est représentable car $M'/A$ existe « universellement »
et alors $M = M'/A$ représente $F$. \ill{} \ill{} \ill{} \ill{} \ill{} $A$
est \emph{propre} et plat sur $S$.

Or $F'(T)$ s'explicite comme \ill{} sous-ensemble de
$\underline{\mathrm{Hilb}}_{A/S}(T)$, \ill{} \ill{} vérifie également \ill{}
$C$ \emph{d'après} \struck{dans} $A_{T}$, \ill{} \ill{} le \ill{} jacobienne
\ill{} \ill{} naturel $J^{0}_{C/S} \longrightarrow A_{T}$ \ill{} \ill{}
\add{\ill{} déduisant} \emph{Matsusaka}, \struck{Par} \ill{} exiger \ill{}
\ill{} \ill{} plongement $C \to A_{T}$, \struck{si $C \to A_{T}$ soit} \ill{}
$C$ \ill{} de dernière. D'après le \emph{théorème de Hilbert}, \ill{} $\Theta$
\ill{} absolument $J^{0}_{C/S} \to A_{T}$ \ill{}

\page{8}

\struck{\ill{}}

\ill{} d'après ceci ? (\emph{Une déformation infinitésimale} \add{du premier
ordre} \emph{de $C$ dans $A$} [si $C$ est \ill{} \ill{} \emph{positionnée} dans
$A$] \add{ne} \struck{\ill{}} \emph{translation infinitésimale} de $A$ (qui est
bien déterminée d'ailleurs, puisque \ill{} pas d'automorphismes
\add{infinitésimaux}).

Or \struck{comme} les déformations infinitésimales en question sont définies par
un élément de $H^{0}(C, \underline{N})$, où $\underline{N}$ est le faisceau
normal de $C$ \struck{dans} $A$.

Considérant la suite exacte
\[
  0 \longrightarrow \underline{T} \longrightarrow \underline{V}
    \longrightarrow \underline{N} \longrightarrow 0
\]
où $\underline{V}$ est induit par le faisceau tangent à $A$ (un \ill{}
constant de rang $g$, $\check{\Omega}$), $\underline{T}$ est le faisceau
tangent \ill{} : $C$ \ill{} de degré $2(1-g) < 0$, \struck{et} \ill{}
\[
  H^{0}(C, \underline{T}) = 0, \qquad
  H^{0}(C, \underline{V}) \simeq k^{g} \simeq
  H^{0}(A, \check{\Omega}_{A/k}),
\]
et la question \ill{} donc \ill{} $H^{0}(C, \underline{V}) \to H^{0}(C,
\underline{N})$ \ill{} \emph{surjectif}, i.e.
\[
  H^{1}(C, \underline{T}) \longrightarrow H^{1}(C, \underline{V})
  \simeq \struck{\ill{}} H^{1}(A, \check{\Omega}_{A/k})
\]
\emph{injectif}, i.e.

\page{9}

\textbf{Corollaire}\quad \emph{Toute déformation infinitésimale} \add{du
premier ordre} \emph{de la} \emph{structure de $C$} qui induit une déformation
infinitésimale \emph{triviale} de la \emph{structure} de $J^{0}_{C/S}$,
\emph{est triviale}.

Traduisant par dualité, on est ramené à la question si
\[
  H^{0}(C, \check{\underline{V}} \otimes \underline{\Omega})
  = H^{0}(C, \underline{\Omega}) \otimes H^{0}(C, \underline{\Omega})
  \longrightarrow H^{0}(C, \underline{\Omega}^{\otimes 2})
\]
(où $\Omega = \Omega^{1}_{C/k}$) est \emph{surjectif}. Or ce dernier fait est
vrai \emph{ssi} $C$ est non \emph{hyperelliptique} ; dans le cas
\struck{\ill{}} hyperelliptique, on sait par \ill{} \ill{} \ill{} \ill{}
\ill{} \ill{} $H^{0}(\underline{V}) \to H^{0}(\underline{N})$ \ill{} de
dimension 1, donc \ill{} \ill{} $S$) la fibre de $M$ en $s$ \ill{}
$\mathrm{Spec}\,k(s)$.

(ii)\quad On voit \ill{} que $M \to S$ est \ill{} \ill{} \ill{} \ill{} \ill{}
\emph{relatif}) \ill{} \ill{} \ill{} \ill{} :

\textbf{Corollaire}\quad \struck{Supposons} \ill{} $S = \mathrm{Spec}(k)$,
$\underline{V}$ \ill{} \ill{} relativement discrète, \struck{Soit} $y$ \ill{}
\ill{} \ill{} \ill{} le \ill{} $M \to \underline{V}$ \ill{} i.e.\ que le \ill{}
de $A$ sur $\underline{V}$ \ill{} des jacobiennes. \struck{Alors} Soit $C_{0}$
une courbe sur \ill{}

\note{le bord gauche de la page porte, en biais, une note de sa main : on y lit
« \ill{} finalement \ill{} », et rien de sûr au-delà}

\page{10}

\ill{} (compatible \ill{} \ill{}),
\[
  J^{0}_{C_{0}/k(s)} \xrightarrow{\;\sim\;} A_{s}.
\]
\ill{} courbe $C$ sur $S$, \struck{$C$} dans \ill{} plus \ill{} isom.\ de
$J^{0}_{C/S}$ \ill{} $A$.

\ill{}, la courbe $C_{0}$ \ill{} \uncertain{sera} évident \ill{} \ill{} par
$J^{0}_{C/S} \to A$ de façon unique \ill{}

\note{les deux tiers inférieurs de la page sont un bloc encadré puis annulé par
une nappe de diagonales, par-dessus une écriture déjà serrée : il n'est pas
lisible et n'est pas donné}

\page{11}

Ceci utilisant \ill{} le fait que $A$ \ill{} (\ill{} une jacobienne) nous dit
que \struck{cette courbe} $C_{a}$ \ill{} a bien une \emph{multiplicité} 1, et
que $J^{0}_{C_{a}/k(a)}$ \ill{} \ill{} produit des jacobiennes \ill{} des
composantes de $C_{a}$.] (« \ill{} »)

Mais un résultat de \uncertain{Matsusaka} nous apprend que si une \ill{}
\ill{} $A$ un \emph{diviseur irréductible} \add{positif} $\Theta_{a}$ \ill{}
\struck{\ill{}} tel que
\[
  \tfrac{1}{n!}\,\deg \Theta_{a}^{n} = 1,
  \qquad
  \Theta_{a}^{\,n-1} \equiv (n-1)!\, Z
\]
\ill{} \ill{} \ill{} $Z$ \struck{soit} \ill{} \ill{} une \emph{courbe} \ill{}
\emph{positif}, alors \ill{} \ill{} \emph{irréductible}, [et de plus \ill{}
\ill{} \ill{} \ill{} $A$ \ill{} bien une \ill{} \ill{} : sa \emph{jacobienne}]
\ill{} $A$, \ill{} $\Theta$ définit \ill{} l'immersion \ill{} la polarisation
\ill{}

\struck{Que l'hypothèse \ill{}}

\ill{} $\Theta_{y}$ \ill{} correspondantes \ill{} $\Theta_{a}$, $\Theta$, dire
\ill{} que
\[
  \deg(\Theta_{y}^{n})/n! = 1
  \quad\text{et}\quad
  \deg(\Theta_{a}^{n}/n!) = 1,
\]
\struck{\ill{} $C_{a}$ était irréd.} et
\[
  \Theta_{y}^{\,n-1} \equiv \struck{\ill{}}\ (n-1)!\, C_{y},
  \qquad
  \Theta_{a}^{\,n-1} \equiv (n-1)!\, C_{a}.
\]
\struck{irréductible, $\Theta_{a}$ \ill{}}

\note{la formule est écrite deux fois, la première biffée ; seule la seconde
est donnée. La classe que l'énoncé général appelle $Z$ est ici la courbe
elle-même, $C_{y}$ ou $C_{a}$}

\page{12}

\emph{irréductible au sens du th.\ de \emph{Matsusaka}}.

\struck{\ill{}} \add{\ill{}} \ill{} \ill{} $A$ \ill{} une jacobienne, \ill{}
\ill{} \struck{\ill{}} \ill{} de \emph{multiplicité} \ill{}

Matsusaka affirme qu'il \ill{} \ill{} mais \ill{} \ill{} dit. $D$ positif
\ill{} il existe \struck{\ill{}} \ill{} \ill{} \emph{irréductible} \ill{}
\ill{}) une $A$ \ill{} \ill{} : \ill{} \uncertain{jacobiennes} \ill{}

\note{le haut de la page est traversé de deux longues diagonales et le peu qui
se lit est donné tel quel}

\ill{} une \ill{}, et \ill{} \ill{} dans le cas de \uncertain{Matsusaka}
\ill{} ; \struck{donc} $C_{a}$ est \emph{irréductible} \ill{} sur $k$. Donc
$C$ \ill{} \emph{et il} \ill{} \ill{} sur $S$, maintenant \ill{}. \ill{}
donnent d'ailleurs ceci :

La démonstration des \ill{}

\textbf{Corollaire 1}\quad Soit $T \cdot \struck{\ill{}} = \overline{f(M)}$, où
$f$ \ill{} le morphisme structural $M \to S$. Pour que $t \in T$ \ill{} de
$f(M)$, il f.\ et s.\ que le diviseur \ill{} \ill{} la polarisation de $A_{t}$
[\ill{} à translation près \ill{} \add{nulle,}] $\deg \Theta^{g}/g! = 1$] soit
\emph{irréductible}.

\page{13}

\struck{\textbf{Corollaire 2}\quad Soient $C$, $C'$ deux courbes sur $S$, munies
d'isomorphismes \ill{} relatives \ill{} $\varphi : J^{0}_{C/S} \to A$ et
$\varphi' : J^{0}_{C'/S} \to A$. Alors la \ill{}}

\note{le premier bloc de la page est encadré puis annulé d'un long trait
sinueux ; il est donné tel qu'il se lit}

\textbf{Corollaire 3}\quad (i)\quad Supposons que $S$ soit \add{réduits, et}
\ill{} \ill{} \ill{} $s \in S$ tel que $A_{s}$ soit \ill{}
\emph{hyperelliptique} \add{polarisé}. Pour que \ill{} une \emph{jacobienne}
\add{polarisée}, il f.\ et s.\ que $A$ ait une \ill{} \ill{} \ill{} soient ;
\ill{} \ill{} que les composantes \ill{} soient \ill{} (composantes de $S$)
\ill{} le soient, et que les \ill{} le diviseur de \uncertain{Matsusaka} \ill{}
\ill{} puissent être \emph{irréductibles}.

(ii)\quad Supposons que $S$ soit \emph{réduit} \ill{} aux \struck{\ill{}} $s
\in S$ \ill{} que $A_{s}$ soit \emph{hyperelliptique}. \struck{deux} Soient
$C$, $C'$ deux courbes sur $S$, munies d'isom.\ de \ill{} \ill{} relations
polarisées
\[
  \varphi : J^{0}_{C/S} \longrightarrow A,
  \qquad
  \varphi' : J^{0}_{C'/S} \longrightarrow A.
\]
Alors il existe un $S$-isom.\ \ill{} $u : C \to C'$ et un \ill{} seul \ill{}
compatible \ill{} avec $\varphi$ et $\varphi'$.

\note{le feuillet s'arrête là, sur l'énoncé et sans démonstration}

\end{document}
